% CVPR 2026 Paper Template
\documentclass[10pt,twocolumn,letterpaper]{article}

\usepackage[review]{cvpr}
\usepackage{times}
\usepackage{epsfig}
\usepackage{graphicx}
\usepackage{amsmath}
\usepackage{amssymb}
\usepackage{booktabs}
\usepackage{array}
\usepackage{multirow}

\definecolor{cvprblue}{rgb}{0.21,0.49,0.74}
\usepackage[pagebackref,breaklinks,colorlinks,allcolors=cvprblue]{hyperref}

\def\paperID{*****}
\def\confName{COMP4471}
\def\confYear{2026}

\title{Decoupled Image-to-Video Generation via Latent Flow:\\An Investigation into Embedding-Conditioned Rendering}

\author{ZhuFeiHao\\
The Hong Kong University of Science and Technology\\
{\tt\small fhzhu@connect.ust.hk}
}

\begin{document}
\maketitle

%%%%%%%%% ABSTRACT
\begin{abstract}
Image-to-video generation from a single static image and an action label
requires simultaneously understanding both what is in the scene and how it should
evolve over time. We investigate a decoupled framework inspired by the
Joint-Embedding Predictive Architecture (JEPA): a causal Transformer (Flow Model)
predicts the trajectory of CLIP embeddings in a compact 512-dimensional semantic
space, and a conditioned UNet (Renderer) converts each predicted embedding back
to a pixel frame. We conduct a systematic set of experiments across four Renderer
architectures---FiLM conditioning, cross-attention to CLIP patch tokens,
contrastive conditioning loss, and optical flow prediction---and find that each
variant collapses to a mode where the embedding condition is effectively ignored
and the output is a near-identity transformation of the input frame. We present a
detailed failure analysis identifying the root causes: vanishing embedding motion
signal between adjacent frames, a training objective fundamentally misaligned
with conditioning sensitivity, and the inherent difficulty of encoding spatial
motion into a single global feature vector. Despite the negative result, our
analysis highlights key design challenges for embedding-conditioned video
synthesis and suggests concrete alternative approaches.
\end{abstract}

%%%%%%%%% INTRODUCTION
\section{Introduction}

Image-to-video (I2V) generation aims to produce a temporally coherent video
sequence from a single static frame. When conditioned on an action label
(e.g., ``walking''), the generated video should depict natural motion consistent
with that action. This task is relevant to content creation, video editing, and
synthetic data augmentation for action recognition systems.

Recent large-scale diffusion models such as Stable Video Diffusion and
Wan2.2 have demonstrated impressive results by inflating image diffusion
architectures into the temporal domain. However, these end-to-end approaches
incur substantial computational costs and remain prone to motion flickering
and identity drift.

Inspired by the JEPA framework, we investigate an alternative: decoupling
the problem into a \textbf{Flow Model} that predicts motion in a semantic
embedding space and a \textbf{Renderer} that converts semantic predictions
back to pixels. The key hypothesis is that motion dynamics are easier to
model in a compact, abstract embedding space than in raw pixels. In an ideal
scenario, the Flow Model would learn high-level action dynamics (``a walking
person alternates legs''), and the Renderer would faithfully translate these
semantic blueprints into frames.

We implement this framework using frozen CLIP ViT-B/32 encoders, a causal
Transformer for embedding trajectory prediction, and a UNet-based Renderer
with conditioning injected via FiLM, cross-attention, or explicit flow
prediction heads. We systematically ablate the conditioning mechanism and
loss function across four Renderer variants.

\textbf{Our key finding is negative but instructive}: across all four variants,
the Renderer learns to ignore the embedding condition and produces
near-identical outputs regardless of input conditioning. The root cause is
not any single architectural choice but a fundamental tension between (1) the
minimal variation in CLIP embeddings between adjacent video frames, and (2)
loss functions that penalize deviation from a nearly-static ground truth.

We present a detailed failure analysis quantifying the embedding signal
strength, the collapse behavior across architectures, and the training
dynamics. We believe these findings provide useful guidance for future work
on embedding-conditioned generative models.

%%%%%%%%% RELATED WORK
\section{Related Work}

\subsection{Image-to-Video Generation}
Early video generation methods employed recurrent networks or generative
adversarial networks. More recently, diffusion-based approaches have become
standard. Stable Video Diffusion fine-tunes a pre-trained image diffusion
model on video data by adding temporal layers. Wan2.2 follows a similar
approach at larger scale. These models generally operate in a VAE latent
space, simultaneously modeling spatial structure and temporal dynamics in one
pipeline. We experiment with a fundamentally different architecture that
attempts to factor these responsibilities across two modules.

\subsection{Latent Predictive Architectures}
The JEPA framework proposes that predictive learning should operate in
representation space rather than input space. V-JEPA applies this
principle to self-supervised video learning by predicting masked
spatio-temporal features. Our Flow Model follows the same intuition: predict
the trajectory of CLIP embeddings rather than raw pixels. Where our work
differs is the downstream rendering step---converting predicted embeddings
back to pixels---which proves to be the critical bottleneck.

\subsection{Conditioned Image Generation}
Feature-wise Linear Modulation (FiLM) and cross-attention are
standard techniques for conditioning image generation on external
information. FiLM applies channel-wise affine transformations
$\gamma(c) \odot x + \beta(c)$ predicted from a conditioning vector $c$,
and has been widely used in visual question answering, style transfer, and
video prediction. Cross-attention, popularized by the Transformer
architecture, allows each spatial position in a feature map to attend to a
set of conditioning tokens. Our experiments test both mechanisms
independently and in combination, and find that neither is sufficient when
the conditioning signal (adjacent-frame CLIP embeddings) varies by less
than 3\% in cosine distance.

\subsection{Optical Flow and Warping}
Frame interpolation and video prediction often employ explicit optical flow
prediction followed by warping. Methods such as RAFT achieve
state-of-the-art optical flow estimation through iterative refinement,
while simpler approaches use Farneback or TV-L1 flow as a supervision
signal. Our final Renderer variant adopts a flow-based approach: a UNet
backbone predicts a dense flow field conditioned on CLIP embeddings, and
the previous frame is warped via differentiable bilinear sampling. We use
Farneback optical flow as pseudo-supervision and apply total variation
regularization to encourage spatial smoothness.

\subsection{Action-Conditioned Video Generation}
Motion generation conditioned on action labels has been studied in
controlled settings such as video game environments and simplified
human motion benchmarks. Text-to-video models such as Make-A-Video and
Imagen Video use large language model embeddings for open-vocabulary
conditioning at scale. Our setting is intermediate: we use closed-set
action labels encoded through CLIP text embeddings, providing a clean
semantic signal in the same embedding space as the visual features.

%%%%%%%%% DATA
\section{Data}

We use the HMDB51 dataset, containing 6,754 video clips across 51
human action categories such as ``walking,'' ``drinking,'' ``brushing hair,''
and ``climbing stairs.'' Videos average approximately 3 seconds in duration at
25--30 fps. The dataset is split into training (70\%) and test (30\%) sets
following the first official split.

\subsection{Preprocessing}
For each video, we sample 17 consecutive frames ($n=16$ prediction steps) at
a stride of 2, and resize all frames to $224\times 224$ pixels. CLIP
embeddings (CLS token + 49 patch tokens) are pre-extracted for the entire
dataset (242MB) to avoid redundant forward passes during Flow Model training.
The Renderer samples random frame pairs $(I_{t-1}, I_t)$ from each video
with a gap of 1 frame.

\subsection{Embedding Analysis}
A key diagnostic we perform early is measuring the variation in CLIP
embeddings across frames. Table~\ref{tab:embed_analysis} summarizes pairwise
cosine similarities.

\begin{table}[h]
\centering
\caption{CLIP ViT-B/32 embedding cosine similarity. Same-video frames are
highly similar ($\text{1-cos}=0.024$ for adjacent frames), limiting the
information available for frame-to-frame motion prediction.}
\label{tab:embed_analysis}
\begin{tabular}{lr}
\toprule
\textbf{Comparison} & \textbf{Mean cos} \\
\midrule
Same video, distant frames ($\Delta t \approx 60$) & 0.862 \\
Same video, adjacent frames ($\Delta t = 1$)       & 0.976 \\
Different videos, same action class                & 0.664 \\
Different action classes                           & 0.564 \\
\bottomrule
\end{tabular}
\end{table}

The critical observation is that adjacent-frame embeddings differ by only
$\text{1-cos} \approx 0.024$ on average. This means the embedding trajectory
that the Flow Model must predict, and that the Renderer must use as a
conditioning signal, contains very little per-frame variation. Any
learning algorithm must amplify this tiny signal by orders of magnitude
to produce meaningful pixel-space changes.

%%%%%%%%% METHODS
\section{Methods}

Our framework consists of three frozen encoders and two trainable modules.
Figure~\ref{fig:overview} illustrates the overall pipeline.

\begin{figure*}[t]
\centering
\fbox{\rule{0.9\textwidth}{0pt}\rule{0pt}{2.5cm}}
\caption{Overall pipeline. Frozen CLIP encoders produce image embeddings
$E_t$ and label embeddings $L$. The Flow Model (causal Transformer) predicts
$(\hat{E}_1, \dots, \hat{E}_n)$ from $(L, E_0)$. The Renderer (UNet with
various conditioning mechanisms) produces each frame $I_t$ from $I_{t-1}$
conditioned on $\hat{E}_t$ and $L$.}
\label{fig:overview}
\end{figure*}

\subsection{Encoders}

\textbf{Image Encoder.} Frozen CLIP ViT-B/32. We extract both the 512-dim
CLS embedding $E_t$ and the 49 spatial patch tokens ($7\times 7$ grid,
512-dim each) from the final vision transformer layer.

\textbf{Label Encoder.} Frozen CLIP text encoder. Action labels are
converted from directory names to natural text (e.g., ``brush hair'') and
encoded to 512-dim semantic embeddings $L$.

\subsection{Flow Model}

A 4-layer causal Transformer encoder (12.9M parameters) with 8 attention
heads and feedforward dimension 2048. The input sequence is
$[L, E_0, E_1, \dots, E_{n-1}]$, and the model is trained with shifted
teacher forcing: position $t$ predicts $\hat{E}_{t}$ from
$(L, E_{0, \dots, t-1})$.

We compare two loss functions:
\begin{itemize}
    \item \textbf{MSE}: $\mathcal{L} = \frac{1}{n}\sum_t \|\hat{E}_t - E_t\|_2^2$
    \item \textbf{Cosine}: $\mathcal{L} = \frac{1}{n}\sum_t (1 - \cos(\hat{E}_t, E_t))$
\end{itemize}

\subsection{Renderer}

The Renderer is a UNet (4 down/up blocks, 40--43M parameters depending on
variant) that takes $I_{t-1}$ as input and produces an output conditioned
on $\hat{E}_t$, $L$, and patch tokens. We experiment with four variants:

\begin{enumerate}
    \item \textbf{FiLM (baseline)}: channel-wise Feature-wise Linear
    Modulation from a condition vector $[\hat{E}_t; L] \in \mathbb{R}^{1024}$,
    applied at each residual block. Output: pixel residual $\Delta I_t$.
    Loss: $\text{L1}(\Delta I_t, I_t - I_{t-1})$.

    \item \textbf{FiLM + Cross-Attention}: FiLM conditioning plus
    multi-head cross-attention from UNet feature maps to CLIP patch tokens
    at the bottleneck and first two upsampling blocks. Output: pixel
    residual $\Delta I_t$. Loss: L1.

    \item \textbf{FiLM + Cross-Attention + Contrastive}: same as variant 2,
    plus a contrastive conditioning loss: for a batch, we compute
    $\mathcal{L}_{\text{cont}} = \|\cos(\Delta I_{\text{real}}, \Delta I_{\text{fake}}) - \cos(E_{\text{real}}, E_{\text{fake}})\|_2^2$,
    where fake conditions are shuffled embeddings from other videos in the
    batch. This explicitly penalizes the Renderer for producing identical
    outputs from different conditions.

    \item \textbf{Flow Prediction}: same UNet backbone with FiLM +
    cross-attention, but the output head predicts a dense optical flow field
    $\in \mathbb{R}^{2 \times 224 \times 224}$. The previous frame is warped
    via the predicted flow: $I_t = \text{warp}(I_{t-1}, \text{flow})$.
    Loss: $\text{L1}(I_t^{\text{warped}}, I_t^{\text{gt}}) + \lambda_{\text{TV}} \cdot \text{TV}(\text{flow}) + \lambda_{\text{FB}} \cdot \text{L1}(\text{flow}, \text{Farneback})$.
\end{enumerate}

\subsection{Training Details}

All models are trained with AdamW ($lr=10^{-4}$, $\beta_1=0.9$,
$\beta_2=0.999$, weight decay $10^{-2}$) and a cosine annealing learning
rate schedule over 100 epochs. We use a batch size of 8 with gradient
accumulation over 2 steps (effective batch size 16). Gradient clipping
with maximum norm 1.0 is applied. All experiments were conducted on a
single NVIDIA RTX 4060 Laptop GPU (8GB VRAM). Complete training of the Flow
Model takes approximately 4 hours; each Renderer variant takes
approximately 7 hours.

\textbf{Flow Model training data pipeline.} Because passing CUDA models
through PyTorch DataLoader workers is not supported on Windows, we
pre-extract CLIP embeddings for all 6,754 videos (17 frames each) and save
them to disk as a single 242MB file. The Flow Model training loop loads
pre-extracted embeddings directly, requiring no CLIP forward passes during
training. This eliminates the DataLoader CUDA issue and accelerates
training by approximately 3$\times$.

\textbf{Renderer training data pipeline.} The Renderer uses
HMDB51FrameDataset, which loads raw frames from disk (via OpenCV) and
performs image transformations in the DataLoader workers. CLIP encoding
happens in the main training loop where CUDA is available. Farneback
optical flow is computed on CPU using OpenCV, adding approximately 30\%
overhead per batch.

\subsection{Implementation Notes}

\textbf{Shifted teacher forcing.} An earlier version of the Flow Model used
``zero-placeholder'' inputs where the model received $[L, E_0, \mathbf{0},
\dots, \mathbf{0}]$ and was expected to predict all future embeddings in a
single forward pass. This created a train-inference mismatch: at inference
time, the model sees its own predictions rather than zeros. Shifted teacher
forcing, where the model receives the real embedding sequence and predicts
one step ahead at each position, aligns training and inference exactly. We
found this improved embedding trajectory quality significantly (from 2.0\%
to 82.5\% of GT motion, combined with the cosine loss change).

\textbf{Multi-scale flow prediction.} In the flow prediction Renderer
variant, we considered predicting flow at multiple decoder resolutions and
combining them via coarse-to-fine refinement (similar to RAFT). However,
initial experiments showed that even a single-scale flow head struggled to
produce consistent non-zero outputs, so we deferred multi-scale
investigation to future work.

%%%%%%%%% EXPERIMENTS
\section{Experiments}

\subsection{Evaluation Protocol}

We evaluate using the following metrics:
\begin{itemize}
    \item \textbf{Embedding Trajectory MSE}: MSE between predicted and
    ground-truth CLIP embedding sequences.
    \item \textbf{Inter-frame 1-cos Distance}: $1 - \cos(\hat{E}_t, \hat{E}_{t+1})$,
    measuring how much ``motion'' the predicted embeddings contain.
    \item \textbf{Conditioning Sensitivity}: $\frac{\|\text{Renderer}(I, E_t) - \text{Renderer}(I, E_0)\|_1}{\|\text{Renderer}(I, E_0)\|_1 + \epsilon}$,
    measuring whether changing the embedding changes the Renderer output.
\end{itemize}

\subsection{Flow Model Results}

Table~\ref{tab:flow} summarizes the Flow Model results.

\begin{table}[h]
\centering
\caption{Flow Model performance with MSE vs. cosine similarity loss.
Cosine loss prevents embedding collapse while maintaining prediction
accuracy.}
\label{tab:flow}
\begin{tabular}{lcc}
\toprule
\textbf{Metric} & \textbf{MSE Loss} & \textbf{Cosine Loss} \\
\midrule
Train loss (final)         & 0.0116  & 0.0264  \\
Inter-frame 1-cos (GT)     & 0.0290  & 0.0290  \\
Inter-frame 1-cos (pred)   & 0.0006  & 0.0239  \\
\% of GT motion preserved  & 2.0\%   & 82.5\%  \\
\bottomrule
\end{tabular}
\end{table}

The MSE-trained Flow Model achieves low MSE (0.0116) but collapses to
predicting nearly identical embeddings at every timestep (1-cos $\approx
0.0006$). Cosine loss resolves this: predicted embeddings retain 82.5\% of
the ground-truth motion variation while maintaining comparable embedding
quality (MSE between predicted and GT embeddings actually improves in terms
of cosine distance).

\subsection{Renderer Results}

Table~\ref{tab:renderer} presents the core negative result: across all four
architectures, the Renderer learns to ignore the embedding condition.

\begin{table}[h]
\centering
\caption{Renderer conditioning sensitivity across four architecture
variants. A ratio of $\sim\!1.0$ indicates the model produces identical
outputs regardless of the embedding condition.}
\label{tab:renderer}
\begin{tabular}{lc}
\toprule
\textbf{Renderer Variant} & \textbf{Sensitivity Ratio} \\
\midrule
1. FiLM + L1                                  & 1.03x \\
2. FiLM + Cross-Attn + L1                     & 0.98x \\
3. FiLM + Cross-Attn + L1 + Contrastive       & 0.98x \\
4. FiLM + Cross-Attn + Flow Pred + Rec + TV + FB & 1.01x \\
\bottomrule
\end{tabular}
\end{table}

For all variants, the conditioning sensitivity ratio is approximately 1.0,
meaning the magnitude of the Renderer's output does not meaningfully change
when the conditioning embedding changes. In variants 1--3 (residual
prediction), the absolute output magnitude converges to $\approx 0.005$---essentially
near-zero. In variant 4 (flow prediction), the predicted flow field
converges back to zero after initial non-zero activity from random
initialization.

\subsection{Diagnostic Experiments}

To isolate the failure mode, we conduct several diagnostic tests:

\textbf{Embedding distance test.} We feed the Renderer embeddings from
frames at varying temporal distances ($[t, t+1, t+4, t+8, t+16]$) and
measure the output magnitude. Even with frame $t+16$ (1-cos $=0.31$, a
meaningful change), the output magnitude remains unchanged from the $t+1$
case (1-cos $=0.03$). This confirms that the model is not merely limited
by the weak signal in adjacent frames---it has learned to ignore the
embedding entirely.

\textbf{Farneback flow as pseudo-target.} We compute Farneback optical flow
between adjacent frames (mean magnitude 0.39 px, max 2.84 px) and use it
as a supervised target. The flow head initially produces non-zero flow but
converges back to zero as the reconstruction and smoothness losses dominate.

\textbf{CLIP feature discrimination.} We verify that CLIP embeddings
\emph{do} contain frame-to-frame information: the cosine similarity gap
between same-video adjacent frames (0.976) and different videos of the
same class (0.664) is 0.312---a substantial difference. The information
exists, but translating it from a 512-dim global vector into spatially
meaningful pixel changes proves fundamentally difficult with our
architectures.

%%%%%%%%% DISCUSSION
\section{Discussion: Why Did It Fail?}

We identify four interrelated factors contributing to the Renderer's
failure:

\subsection{Factor 1: Vanishing Embedding Gradient}
Adjacent-frame CLIP embeddings differ by 1-cos $\approx 0.024$. For the
Flow Model, this is a learnable signal (we achieve 82.5\% of GT variation).
But for the Renderer, this means the conditioning vector $[E_t; L]$ changes
by $<3$\% between frames---far too little for the UNet to learn a
meaningful mapping from condition to output.

\subsection{Factor 2: Loss Function Misalignment}
All our loss functions (L1 on residuals, L1 on reconstruction after
warping) have the property that the \emph{optimal constant prediction} is
already near the ground truth: $\Delta I_t \approx 0$ and $\text{flow}
\approx 0$. The model can achieve a respectably low loss by ignoring the
condition and outputting a near-identity transformation. The contrastive
conditioning loss in variant 3 was designed to address this, but it only
constrains the \emph{direction} of output variation, not its magnitude.
A constant-near-zero output with slightly different noise patterns per
condition can satisfy the contrastive constraint while still producing no
meaningful motion.

\subsection{Factor 3: FiLM Cannot Express Spatial Motion}
Feature-wise Linear Modulation applies channel-wise scaling and shifting
that is uniform across all spatial positions. This can adjust global
brightness or contrast but cannot communicate ``the hand should move to the
upper-right while the background stays still.'' Cross-attention to CLIP
patch tokens (variant 2) was an attempt to introduce spatial conditioning,
but the patch tokens from CLIP ViT capture \emph{what} is in each region,
not \emph{how it should move}.

\subsection{Factor 4: Representation Mismatch}
CLIP embeddings are trained for image-text alignment; they encode semantic
content (``a person walking'') but not spatial configuration (``the right
leg is bent at 30 degrees and positioned at coordinates (x, y)''). The
information needed to predict pixel-level frame changes requires spatial
representations that CLIP's global pooling discards. Using patch tokens
helps but does not fully resolve this, as CLIP patch tokens are not trained
for dense correspondence tasks.

\subsection{Limitations of Our Study}

We acknowledge several limitations: (1) we only explored CLIP ViT-B/32 as
the embedding backbone---larger ViT variants or dense-prediction models
like DINOv2 may have different properties; (2) we did not attempt joint
end-to-end training of the Flow Model and Renderer, which could allow the
Flow Model to learn embeddings specifically optimized for rendering; (3)
we used a relatively small dataset (6.7k videos), and the collapse
behavior may differ at larger scale where more diverse motion patterns
provide a stronger learning signal; (4) HMDB51 videos are constrained to
specific human action categories, and the learned dynamics may not
generalize to more open-ended video generation.

\subsection{What We Learned}

Despite the negative result, this investigation yielded several lessons
that may benefit future work in this area:

\textbf{Embedding choice is critical.} The choice of CLIP ViT-B/32 as our
embedding backbone was motivated by its joint vision-language space and
compact dimensionality. However, CLIP embeddings are optimized for semantic
discrimination, not spatial correspondence. For a decoupled I2V pipeline
to succeed, the embedding space must encode both \emph{what} is in the
scene and \emph{where} things are spatially---properties that CLIP's
contrastive training objective does not guarantee.

\textbf{Loss functions define the optimization landscape.} We observed
that all tested pixel-space loss functions (L1, L1 + TV, L1 + Farneback)
contain a near-identity trivial solution ($\Delta I \approx 0$, flow
$\approx 0$). This is not specific to our architecture---any model trained
with a pixel reconstruction loss on adjacent-frame pairs faces the same
challenge. Loss functions that operate in a learned feature space
(perceptual loss, adversarial loss) may provide a stronger signal for
motion learning.

\textbf{The JEPA philosophy needs a rendering pathway.} JEPA's core
insight---predict in representation space, not input space---proved
effective for our Flow Model (82.5\% of GT motion). But the companion
challenge of \emph{decoding} those predictions back to the pixel domain is
equally hard and received comparatively little attention in our
implementation. The rendering problem deserves its own careful treatment
rather than being treated as an architectural afterthought.

\subsection{Future Work}

Several promising directions emerge from this analysis.

\textbf{Dense spatial embeddings.} Replacing CLIP with a model trained for
dense visual correspondence---such as DINOv2, Stable Diffusion's VAE
latents, or a self-supervised optical flow network---could provide
embeddings that naturally encode spatial configuration. A trajectory in
DINOv2 patch feature space would carry richer spatial information than a
trajectory in CLIP CLS space.

\textbf{Joint end-to-end training.} Training the Flow Model and Renderer
jointly, with gradients flowing from the rendering loss back through the
Flow Model, would allow the embedding space to adapt to the needs of the
Renderer. The Flow Model would no longer be constrained to match CLIP's
embedding structure and could learn features specifically optimized for
frame rendering.

\textbf{Iterative refinement.} Rather than predicting the full frame
transformation in a single forward pass, a multi-step refinement process
(similar to diffusion models or recurrent refinement) could gradually build
up the output from an initial near-identity estimate. This would allow the
model to learn motion in stages---coarse structure first, then fine
details.

\textbf{Adversarial and perceptual losses.} GAN-based losses or CLIP-based
perceptual losses directly penalize lack of motion and unrealistic static
output, potentially breaking the near-identity local minimum that trapped
all of our L1-trained models.

\textbf{Explicit motion field supervision.} Using a pre-trained optical
flow network (e.g., RAFT) instead of Farneback as a supervision target
could provide a cleaner learning signal. Combined with a rendering loss,
this could guide the model toward physically plausible motion while
maintaining visual fidelity.

%%%%%%%%% CONCLUSION
\section{Conclusion}

We set out to build a decoupled image-to-video generation system in which a
causal Transformer predicts action-conditioned embedding trajectories and a
conditioned UNet renders those trajectories to pixels. The Flow Model
component was partially successful: with cosine similarity loss it learned
to predict embedding trajectories preserving 82.5\% of ground-truth motion
variation, while MSE loss led to complete collapse. However, the Renderer
failed across all four conditioning architectures tested: FiLM,
cross-attention, contrastive conditioning, and optical flow prediction.

Our analysis suggests the failure is not due to any single architectural
choice but stems from a fundamental mismatch: the Renderer is asked to
extract spatially-precise motion instructions from a conditioning signal
that varies by less than 3\% between frames and that was never designed to
encode spatial configuration. All tested loss functions contain a
near-identity trivial solution that dominates the optimization.

\textbf{Future directions} suggested by this analysis include: (1)
replacing the global CLIP embedding with dense spatial features from a
model trained for correspondence (e.g., DINOv2, optical flow networks);
(2) training the Flow Model and Renderer jointly end-to-end, so the
Flow Model learns to produce embeddings that are directly useful for
rendering rather than matching a pre-trained representation; (3) replacing
the single-step rendering approach with an iterative or diffusion-based
process where the model can refine its output over multiple steps; (4)
using adversarial or perceptual losses that penalize the absence of motion
rather than L1 deviation from a static ground truth.

{
    \small
    \bibliographystyle{ieeenat_fullname}
    \begin{thebibliography}{99}
        \bibitem{blat2023} A.~Blattmann, T.~Dockhorn, S.~Kulal, et al.
        ``Stable Video Diffusion: Scaling Latent Video Diffusion Models to
        Large Datasets.'' \textit{arXiv:2311.15127}, 2023.
        \bibitem{wan2026} Alibaba. ``Wan2.2: Large-Scale Image-to-Video
        Generation.'' 2025--2026.
        \bibitem{lecun2022} Y.~LeCun. ``A Path Towards Autonomous Machine
        Intelligence.'' \textit{OpenReview}, 2022.
        \bibitem{bardes2024} A.~Bardes, Q.~Garrido, J.~Ponce, et al.
        ``Revisiting Feature Prediction for Learning Visual Representations
        from Video.'' \textit{arXiv:2404.08471}, 2024.
        \bibitem{radford2021} A.~Radford, J.~W.~Kim, C.~Hallacy, et al.
        ``Learning Transferable Visual Models From Natural Language
        Supervision.'' \textit{ICML}, 2021.
        \bibitem{kuehne2011} H.~Kuehne, H.~Jhuang, E.~Garrote, et al.
        ``HMDB: A Large Video Database for Human Motion Recognition.''
        \textit{ICCV}, 2011.
        \bibitem{perez2018} E.~Perez, F.~Strub, H.~de Vries, et al.
        ``FiLM: Visual Reasoning with a General Conditioning Layer.''
        \textit{AAAI}, 2018.
        \bibitem{zhang2018} R.~Zhang, P.~Isola, A.~A.~Efros, et al.
        ``The Unreasonable Effectiveness of Deep Features as a Perceptual
        Metric.'' \textit{CVPR}, 2018.
        \bibitem{vaswani2017} A.~Vaswani, N.~Shazeer, N.~Parmar, et al.
        ``Attention Is All You Need.'' \textit{NeurIPS}, 2017.
        \bibitem{isola2017} P.~Isola, J.-Y.~Zhu, T.~Zhou, et al.
        ``Image-to-Image Translation with Conditional Adversarial
        Networks.'' \textit{CVPR}, 2017.
        \bibitem{ronneberger2015} O.~Ronneberger, P.~Fischer, T.~Brox.
        ``U-Net: Convolutional Networks for Biomedical Image Segmentation.''
        \textit{MICCAI}, 2015.
        \bibitem{oh2015} J.~Oh, X.~Guo, H.~Lee, et al.
        ``Action-Conditional Video Prediction using Deep Networks in Atari
        Games.'' \textit{NIPS}, 2015.
        \bibitem{srivastava2015} N.~Srivastava, E.~Mansimov,
        R.~Salakhutdinov. ``Unsupervised Learning of Video Representations
        using LSTMs.'' \textit{ICML}, 2015.
        \bibitem{vondrick2016} C.~Vondrick, H.~Pirsiavash, A.~Torralba.
        ``Generating Videos with Scene Dynamics.'' \textit{NIPS}, 2016.
        \bibitem{singer2022} U.~Singer, A.~Polyak, T.~Hayes, et al.
        ``Make-A-Video: Text-to-Video Generation without Text-Video
        Data.'' \textit{arXiv:2209.14792}, 2022.
        \bibitem{ho2022} J.~Ho, W.~Chan, C.~Saharia, et al.
        ``Imagen Video: High Definition Video Generation with Diffusion
        Models.'' \textit{arXiv:2210.02303}, 2022.
        \bibitem{monfort2019} M.~Monfort, A.~Andonian, B.~Zhou, et al.
        ``Moments in Time Dataset: One Million Videos for Event
        Understanding.'' \textit{TPAMI}, 2019.
    \end{thebibliography}
}

\end{document}
